% Options for packages loaded elsewhere
\PassOptionsToPackage{unicode}{hyperref}
\PassOptionsToPackage{hyphens}{url}
\PassOptionsToPackage{dvipsnames,svgnames,x11names}{xcolor}
%
\documentclass[
  letterpaper,
  DIV=11,
  numbers=noendperiod]{scrartcl}

\usepackage{amsmath,amssymb}
\usepackage{iftex}
\ifPDFTeX
  \usepackage[T1]{fontenc}
  \usepackage[utf8]{inputenc}
  \usepackage{textcomp} % provide euro and other symbols
\else % if luatex or xetex
  \usepackage{unicode-math}
  \defaultfontfeatures{Scale=MatchLowercase}
  \defaultfontfeatures[\rmfamily]{Ligatures=TeX,Scale=1}
\fi
\usepackage{lmodern}
\ifPDFTeX\else  
    % xetex/luatex font selection
\fi
% Use upquote if available, for straight quotes in verbatim environments
\IfFileExists{upquote.sty}{\usepackage{upquote}}{}
\IfFileExists{microtype.sty}{% use microtype if available
  \usepackage[]{microtype}
  \UseMicrotypeSet[protrusion]{basicmath} % disable protrusion for tt fonts
}{}
\makeatletter
\@ifundefined{KOMAClassName}{% if non-KOMA class
  \IfFileExists{parskip.sty}{%
    \usepackage{parskip}
  }{% else
    \setlength{\parindent}{0pt}
    \setlength{\parskip}{6pt plus 2pt minus 1pt}}
}{% if KOMA class
  \KOMAoptions{parskip=half}}
\makeatother
\usepackage{xcolor}
\setlength{\emergencystretch}{3em} % prevent overfull lines
\setcounter{secnumdepth}{-\maxdimen} % remove section numbering
% Make \paragraph and \subparagraph free-standing
\ifx\paragraph\undefined\else
  \let\oldparagraph\paragraph
  \renewcommand{\paragraph}[1]{\oldparagraph{#1}\mbox{}}
\fi
\ifx\subparagraph\undefined\else
  \let\oldsubparagraph\subparagraph
  \renewcommand{\subparagraph}[1]{\oldsubparagraph{#1}\mbox{}}
\fi


\providecommand{\tightlist}{%
  \setlength{\itemsep}{0pt}\setlength{\parskip}{0pt}}\usepackage{longtable,booktabs,array}
\usepackage{calc} % for calculating minipage widths
% Correct order of tables after \paragraph or \subparagraph
\usepackage{etoolbox}
\makeatletter
\patchcmd\longtable{\par}{\if@noskipsec\mbox{}\fi\par}{}{}
\makeatother
% Allow footnotes in longtable head/foot
\IfFileExists{footnotehyper.sty}{\usepackage{footnotehyper}}{\usepackage{footnote}}
\makesavenoteenv{longtable}
\usepackage{graphicx}
\makeatletter
\def\maxwidth{\ifdim\Gin@nat@width>\linewidth\linewidth\else\Gin@nat@width\fi}
\def\maxheight{\ifdim\Gin@nat@height>\textheight\textheight\else\Gin@nat@height\fi}
\makeatother
% Scale images if necessary, so that they will not overflow the page
% margins by default, and it is still possible to overwrite the defaults
% using explicit options in \includegraphics[width, height, ...]{}
\setkeys{Gin}{width=\maxwidth,height=\maxheight,keepaspectratio}
% Set default figure placement to htbp
\makeatletter
\def\fps@figure{htbp}
\makeatother

\KOMAoption{captions}{tableheading}
\makeatletter
\makeatother
\makeatletter
\makeatother
\makeatletter
\@ifpackageloaded{caption}{}{\usepackage{caption}}
\AtBeginDocument{%
\ifdefined\contentsname
  \renewcommand*\contentsname{Table of contents}
\else
  \newcommand\contentsname{Table of contents}
\fi
\ifdefined\listfigurename
  \renewcommand*\listfigurename{List of Figures}
\else
  \newcommand\listfigurename{List of Figures}
\fi
\ifdefined\listtablename
  \renewcommand*\listtablename{List of Tables}
\else
  \newcommand\listtablename{List of Tables}
\fi
\ifdefined\figurename
  \renewcommand*\figurename{Figure}
\else
  \newcommand\figurename{Figure}
\fi
\ifdefined\tablename
  \renewcommand*\tablename{Table}
\else
  \newcommand\tablename{Table}
\fi
}
\@ifpackageloaded{float}{}{\usepackage{float}}
\floatstyle{ruled}
\@ifundefined{c@chapter}{\newfloat{codelisting}{h}{lop}}{\newfloat{codelisting}{h}{lop}[chapter]}
\floatname{codelisting}{Listing}
\newcommand*\listoflistings{\listof{codelisting}{List of Listings}}
\makeatother
\makeatletter
\@ifpackageloaded{caption}{}{\usepackage{caption}}
\@ifpackageloaded{subcaption}{}{\usepackage{subcaption}}
\makeatother
\makeatletter
\@ifpackageloaded{tcolorbox}{}{\usepackage[skins,breakable]{tcolorbox}}
\makeatother
\makeatletter
\@ifundefined{shadecolor}{\definecolor{shadecolor}{rgb}{.97, .97, .97}}
\makeatother
\makeatletter
\makeatother
\makeatletter
\makeatother
\ifLuaTeX
  \usepackage{selnolig}  % disable illegal ligatures
\fi
\IfFileExists{bookmark.sty}{\usepackage{bookmark}}{\usepackage{hyperref}}
\IfFileExists{xurl.sty}{\usepackage{xurl}}{} % add URL line breaks if available
\urlstyle{same} % disable monospaced font for URLs
\hypersetup{
  pdftitle={Ejercicios SOL114},
  pdfauthor={Mauricio Bucca},
  colorlinks=true,
  linkcolor={blue},
  filecolor={Maroon},
  citecolor={Blue},
  urlcolor={Blue},
  pdfcreator={LaTeX via pandoc}}

\title{Ejercicios SOL114}
\author{Mauricio Bucca}
\date{}

\begin{document}
\maketitle
\ifdefined\Shaded\renewenvironment{Shaded}{\begin{tcolorbox}[breakable, interior hidden, enhanced, frame hidden, boxrule=0pt, borderline west={3pt}{0pt}{shadecolor}, sharp corners]}{\end{tcolorbox}}\fi

\begin{enumerate}
\def\labelenumi{\arabic{enumi}.}
\item
  Se lanza una moneda y un dado. Enumera el espacio muestral.
\item
  Proporciona un ejemplo en el que dos eventos sean dependientes.
\item
  Una bolsa contiene 4 bolas rojas, 6 bolas azules y 5 bolas verdes. Si
  se saca una bola al azar, encuentra la probabilidad de que no sea ni
  roja ni azul.
\item
  Si un jugador de baloncesto tiene una precisión de tiros libres del
  90\%, ¿cuál es la probabilidad de que falle dos veces seguidas?
\item
  ¿Cómo se define la independencia de dos eventos en probabilidad?
\item
  Hay 7 bolas rojas y 3 bolas verdes en una caja. Se saca una bola al
  azar, se anota y luego se remueve de la bolsa. Si se saca una bola
  roja primero, ¿cuál es la probabilidad de que una bola verde se saque
  en segundo lugar?
\item
  En una encuesta a 100 personas, 65 beben té, 40 beben café y 10 no
  beben ninguno de los dos. Encuentra la probabilidad de que una persona
  seleccionada al azar beba té y café.
\item
  Un comerciante sabe que el 30\% de sus clientes entran a la tienda
  después de ver un anuncio. De estos, el 10\% realiza una compra. De
  los clientes que no vieron el anuncio, el 5\% realiza una compra. Si
  un cliente realiza una compra, ¿cuál es la probabilidad de que haya
  visto el anuncio?
\item
  En una comunidad, hay dos clubes sociales prominentes. Los miembros
  del Club A tienen un 75\% de probabilidad de asistir a un gran evento
  comunitario, mientras que los miembros del Club B tienen solo un 50\%
  de probabilidad. Dado que el 30\% de la comunidad pertenece al Club A
  y el resto al Club B: si una persona es vista en un evento
  comunitario, ¿cuál es la probabilidad de que sea del Club A?
\item
  Si la probabilidad de A dado B es 0.6 y la probabilidad de B es 0.2,
  ¿cuál es la probabilidad conjunta de A y B?
\item
  Sea \(X\) una variable aleatoria que representa el número de caras
  obtenidas cuando se lanza una moneda 3 veces. Encuentra \(P(X = 2)\).
\item
  Dada una variable aleatoria discreta \(X\) con la siguiente
  distribución de probabilidad: \(P(X = 0) = 0.2\), \(P(X = 1) = 0.3\),
  \(P(X = 2) = 0.5\). Calcula \(\text{E}(X)\).
\item
  Usando el problema anterior, calcula \(\text{Var}(X)\).
\item
  Sean \(X\) y \(Y\) dos variables aleatorias con las siguientes
  propiedades:

  - \(E[X] = 3\), \(\text{Var}(X) = 4\)

  - \(E[Y] = 5\), \(\text{Var}(Y) = 9\)

  Encuentra el valor esperado y la varianza para \(Z = 2X - 3Y\).
\item
  Supongamos que (\(X_1, X_2, \ldots, X_5\)) son variables aleatorias
  independientes e idénticamente distribuidas con ( \$E{[}X\_i{]} = 4\$
  ) y ( \$\text{Var}(X\_i) = 6\$ ).
\end{enumerate}

Encuentra el valor esperado y la varianza de (\$Y = \frac{1}{5}
\sum\_\{i=1\}\^{}\{5\} X\_i\$), el promedio aritmético de las cinco
variables aleatorias.

\begin{enumerate}
\def\labelenumi{\arabic{enumi}.}
\setcounter{enumi}{15}
\tightlist
\item
  Supongamos que tienes 3 variables aleatorias independientes: \(A\),
  \(B\) y \(C\), donde:
\end{enumerate}

\begin{itemize}
\item
  \(E[A] = 3\), \(\text{Var}(A) = 9\)
\item
  \(E[B] = 5\), \(\text{Var}(B) = 4\)
\item
  \(E[C] = 2\), \(\text{Var}(C) = 1\)
\end{itemize}

Calcula el valor esperado y la varianza de \(Z = \frac{A + B + C}{3}\).

\begin{enumerate}
\def\labelenumi{\arabic{enumi}.}
\setcounter{enumi}{16}
\item
  Sea \(Y\) el número de caras obtenidas al lanzar 4 monedas. Encuentra
  la distribución de probabilidad de \(Y\).
\item
  Un juego consiste en sacar una carta de un mazo. Si sacas un as, ganas
  \$50; un rey, ganas \$30; y cualquier otra carta, pierdes \$10.
  Calcula las ganancias esperadas de este juego.
\item
  Si \(Y\) es la suma de dinero ganado después de 10 extracciones,
  encuentra el valor esperado de \(Y\).
\item
  Una librería vende un libro de texto particular con un 30\% de
  descuento con una probabilidad de 0.6 y un 20\% de descuento con una
  probabilidad de 0.4. Si el precio original del libro es de \$50,
  encuentra el precio esperado después del descuento.
\item
  Una empresa fabrica bombillas, el 98\% de las cuales duran más de 1000
  horas y el 2\% falla antes de eso. Si se prueban 50 bombillas, ¿cuál
  es la probabilidad de que al menos 3 fallen antes de las 1000 horas?
\item
  Se sacan tres cartas de un mazo de 52 cartas sin reemplazo. ¿Cuál es
  la probabilidad de que las tres sean ases?
\item
  En una competición de tiro al blanco, un jugador tiene una
  probabilidad del 80\% de acertar el blanco en cada intento. Si dispara
  10 veces, ¿cuál es la probabilidad de que acierte al menos 7 veces
  pero no más de 9?
\item
  Un estudio reveló que el 25\% de los estudiantes bebe café, el 50\%
  bebe té y el 25\% bebe ambos. Si se selecciona un estudiante al azar,
  ¿cuál es la probabilidad de que beba solo té?
\item
  En una ciudad, el 60\% de las personas usa el transporte público y el
  20\% camina al trabajo. Dado que una persona camina al trabajo, hay un
  5\% de probabilidad de que también use el transporte público en algún
  punto. Si una persona fue vista en el transporte público, ¿cuál es la
  probabilidad de que también camine al trabajo?
\item
  Un boxeador tiene un 70\% de probabilidad de ganar un combate si ha
  tenido más de 6 meses de entrenamiento, y un 40\% si ha tenido menos.
  Si el 80\% de los boxeadores han tenido más de 6 meses de
  entrenamiento, y se selecciona un boxeador ganador al azar, ¿cuál es
  la probabilidad de que haya tenido más de 6 meses de entrenamiento?
\item
  Dada una distribución binomial con \(n = 10\) y \(p = 0.5\), calcula
  la probabilidad de que \(X \leq 3\).
\item
  Usando la función de masa de probabilidad binomial, si se lanza una
  moneda 6 veces, ¿cuál es la probabilidad de obtener 2 caras y 4
  sellos?
\item
  Si \(V\) es una variable aleatoria normal estándar, ¿cuál es la
  probabilidad \(P(-2.33 < V < 0)\)?
\item
  En un juego de lanzamiento de dardos, si un jugador acierta el blanco
  central con una probabilidad de 0.1 y el anillo exterior con una
  probabilidad de 0.3, ¿cuál es la varianza de los puntos obtenidos
  después de 10 lanzamientos, considerando que el blanco central da 5
  puntos y el anillo exterior da 2 puntos?
\item
  Si \(U\) es una variable aleatoria que sigue una distribución uniforme
  entre 10 y 20, encuentra \(\text{E}(U)\) y \(\text{Var}(U)\).
\item
  En una caja hay 3 focos defectuosos y 7 buenos. Si se sacan 3 focos al
  azar, ¿cuál es la probabilidad de que al menos uno sea defectuoso?
\item
  Si \(Z\) es una variable aleatoria normal estándar, encuentra el valor
  de \(z\) tal que \(P(Z > z) = 0.05\).
\item
  ¿Cuál es el valor de \(z\) en una distribución \(Z\) tal que
  \(P(Z < z) = 0.10\)?
\item
  Si \(Z\) es una variable aleatoria normal estándar, ¿cuál es la
  probabilidad \(P(-1.96 < Z < 1.96)\)?
\item
  Explica el concepto de una función de densidad de probabilidad (pdf).
\item
  Las alturas de un grupo de atletas están distribuidas uniformemente
  entre 5.5 y 6.5 pies. ¿Cuál es la probabilidad de que un atleta
  seleccionado al azar mida más de 6 pies?
\item
  Un pueblo tiene personas que podrían contraer una enfermedad de un
  vecino enfermo. Cada día, cada persona tiene un 5\% de probabilidad de
  enfermarse. ¿Qué tipo de distribución sigue el número de personas que
  se enferman? Si tenemos 200 personas, en promedio, ¿cuántas se
  enfermarán en un día? ¿Y cuánto puede variar este número? Dado estas
  200 personas, ¿cuál es la probabilidad de que más de 15 de ellas se
  enfermen en un solo día?
\item
  Una región tiene una precipitación anual promedio que puede modelarse
  como una variable aleatoria continua, distribuida normalmente con una
  media de 1000mm y una desviación estándar de 150mm. Encuentra la
  probabilidad de que en un año dado, la precipitación esté entre 800mm
  y 1200mm.
\item
  La altura de los hombres adultos en una cierta población se distribuye
  normalmente con una media de 175 cm y una desviación estándar de 7 cm.
  ¿Qué proporción de esta población mide más de 189 cm?
\item
  Los resultados de un examen en una clase grande se distribuyen
  normalmente con una media de 70 y una desviación estándar de 10. Si la
  puntuación aprobatoria es 50, ¿qué porcentaje de estudiantes reprueba
  el examen?
\item
  En un estudio sobre la relación entre horas de estudio y rendimiento
  académico, se descubrió que el número de horas que los estudiantes
  estudiaban para un examen en particular se distribuía normalmente. El
  tiempo promedio de estudio era de 5 horas con una desviación estándar
  de 0.75 horas. ¿Qué porcentaje de estudiantes estudió entre 4.25 horas
  y 5.75 horas para el examen?
\item
  Un país informa que el número de inmigrantes en un año se distribuye
  normalmente con una media de 5,000 y una desviación estándar de 800.
  También observan que el número de emigrantes (personas que abandonan
  el país) también se distribuye normalmente, con una media de 3,000 y
  una desviación estándar de 500. ¿En qué proporción de años la
  migración neta (inmigrantes - emigrantes) es positiva en más de 1000?
\item
  Sophie tomó un examen estandarizado y obtuvo 86 de 100. Los resultados
  del examen para dos prestigiosas instituciones académicas, Harvard y
  la PUC, donde Sophie está considerando postular, son los siguientes:
  En Harvard, los resultados del examen se distribuyen normalmente con
  una media de 85 y una desviación estándar de 1. En PUC, los resultados
  del examen se distribuyen normalmente con una media de 83 y una
  desviación estándar de 15. Determina la puntuación Z de Sophie para
  ambas instituciones. Basándose en estas puntuaciones Z, ¿en cuál
  institución está Sophie relativamente mejor posicionada con respecto a
  otros solicitantes?
\item
  Supongamos que las puntuaciones de crédito de los solicitantes de
  préstamos en un cierto banco siguen una distribución normal con una
  puntuación promedio de 650 y una desviación estándar de 50. Si el
  banco típicamente aprueba préstamos para solicitantes con puntuaciones
  de crédito en el 20\% superior, ¿cuál es la puntuación mínima de
  crédito requerida para la aprobación de un préstamo?
\item
  En un examen de admisión universitaria, las puntuaciones siguen una
  distribución normal con una puntuación promedio de 1200 y una
  desviación estándar de 100. Si un estudiante obtiene una puntuación de
  1400 en este examen, ¿en qué percentil se encuentra su puntuación
  entre todos los que tomaron el examen?
\item
  Estás investigando la desigualdad de ingresos en una ciudad con una
  población de 50,000 hogares. En esta ciudad, los ingresos están
  distribuidos de manera altamente desigual, con la mayoría de los
  hogares ganando relativamente poco, mientras que unos pocos ganan
  significativamente más. El ingreso promedio de la población es de
  \$50,000, y la varianza es de \$30,000. Supongamos que recopilas datos
  de una muestra aleatoria de 150 hogares de esta población. En el
  contexto de esta distribución de ingresos altamente desigual, aclara
  por qué este intervalo de probabilidad es una herramienta valiosa para
  hacer inferencias sobre la desigualdad de ingresos en toda la ciudad.
\item
  Supongamos que estás estudiando una población con una media de CI
  conocida de 100 y una varianza conocida de 15. Si tomas muestras
  aleatorias de diferentes tamaños (por ejemplo, 30, 75, 150), discute
  cómo se aplica el Teorema del Límite Central (TLC) si el CI sigue una
  distribución uniforme. Incluye reflexiones sobre los efectos del
  tamaño de la muestra.
\item
  Estás llevando a cabo un estudio sobre la distribución de ingresos en
  una ciudad. En la población, se conoce que el ingreso promedio es de
  \$45,000, y la verdadera desviación estándar de la poblaciónes de
  \$8,000. Recopilas una muestra aleatoria de 200 hogares de la
  población. Deriva la distribución de la media muestral y calcula los
  valores que se encuentran a una desviación estándar por debajo y por
  encima del promedio de la media muestral.
\item
  Estás llevando a cabo un estudio sobre la distribución de ingresos en
  una ciudad. En la población, se conoce que el ingreso promedio es de
  \$45,000, y la verdadera desviación estándar de la poblaciónes de
  \$8,000. Recopilas una muestra aleatoria de 10,000 hogares de la
  población. Deriva la distribución de la media muestral y encuentra el
  percentil 10 y el percentil 90 en la distribución de la media
  muestral.
\end{enumerate}

\newpage

\hypertarget{probabilidad-acumulada-en-distribuciuxf3n-normal-standard-mitad-positiva}{%
\subsection{Probabilidad acumulada en distribución Normal Standard
(mitad
positiva)}\label{probabilidad-acumulada-en-distribuciuxf3n-normal-standard-mitad-positiva}}

\begin{longtable}[]{@{}cc@{}}
\toprule\noalign{}
Z=z & \(\Phi(z)\) \\
\midrule\noalign{}
\endhead
\bottomrule\noalign{}
\endlastfoot
0.00 & 0.50000 \\
0.01 & 0.50399 \\
0.02 & 0.50798 \\
0.03 & 0.51197 \\
0.04 & 0.51595 \\
0.05 & 0.51994 \\
0.06 & 0.52392 \\
0.07 & 0.52790 \\
0.08 & 0.53188 \\
0.09 & 0.53586 \\
0.10 & 0.53983 \\
0.11 & 0.54380 \\
0.12 & 0.54776 \\
0.13 & 0.55172 \\
0.14 & 0.55567 \\
0.15 & 0.55962 \\
0.16 & 0.56356 \\
0.17 & 0.56749 \\
0.18 & 0.57142 \\
0.19 & 0.57535 \\
0.20 & 0.57926 \\
0.21 & 0.58317 \\
0.22 & 0.58706 \\
0.23 & 0.59095 \\
0.24 & 0.59483 \\
0.25 & 0.59871 \\
0.26 & 0.60257 \\
0.27 & 0.60642 \\
0.28 & 0.61026 \\
0.29 & 0.61409 \\
0.30 & 0.61791 \\
0.31 & 0.62172 \\
0.32 & 0.62552 \\
0.33 & 0.62930 \\
0.34 & 0.63307 \\
0.35 & 0.63683 \\
0.36 & 0.64058 \\
0.37 & 0.64431 \\
0.38 & 0.64803 \\
0.39 & 0.65173 \\
0.40 & 0.65542 \\
0.41 & 0.65910 \\
0.42 & 0.66276 \\
0.43 & 0.66640 \\
0.44 & 0.67003 \\
0.45 & 0.67364 \\
0.46 & 0.67724 \\
0.47 & 0.68082 \\
0.48 & 0.68439 \\
0.49 & 0.68793 \\
0.50 & 0.69146 \\
0.51 & 0.69497 \\
0.52 & 0.69847 \\
0.53 & 0.70194 \\
0.54 & 0.70540 \\
0.55 & 0.70884 \\
0.56 & 0.71226 \\
0.57 & 0.71566 \\
0.58 & 0.71904 \\
0.59 & 0.72240 \\
0.60 & 0.72575 \\
0.61 & 0.72907 \\
0.62 & 0.73237 \\
0.63 & 0.73565 \\
0.64 & 0.73891 \\
0.65 & 0.74215 \\
0.66 & 0.74537 \\
0.67 & 0.74857 \\
0.68 & 0.75175 \\
0.69 & 0.75490 \\
0.70 & 0.75804 \\
0.71 & 0.76115 \\
0.72 & 0.76424 \\
0.73 & 0.76730 \\
0.74 & 0.77035 \\
0.75 & 0.77337 \\
0.76 & 0.77637 \\
0.77 & 0.77935 \\
0.78 & 0.78230 \\
0.79 & 0.78524 \\
0.80 & 0.78814 \\
0.81 & 0.79103 \\
0.82 & 0.79389 \\
0.83 & 0.79673 \\
0.84 & 0.79955 \\
0.85 & 0.80234 \\
0.86 & 0.80511 \\
0.87 & 0.80785 \\
0.88 & 0.81057 \\
0.89 & 0.81327 \\
0.90 & 0.81594 \\
0.91 & 0.81859 \\
0.92 & 0.82121 \\
0.93 & 0.82381 \\
0.94 & 0.82639 \\
0.95 & 0.82894 \\
0.96 & 0.83147 \\
0.97 & 0.83398 \\
0.98 & 0.83646 \\
0.99 & 0.83891 \\
1.00 & 0.84134 \\
1.01 & 0.84375 \\
1.02 & 0.84614 \\
1.03 & 0.84849 \\
1.04 & 0.85083 \\
1.05 & 0.85314 \\
1.06 & 0.85543 \\
1.07 & 0.85769 \\
1.08 & 0.85993 \\
1.09 & 0.86214 \\
1.10 & 0.86433 \\
1.11 & 0.86650 \\
1.12 & 0.86864 \\
1.13 & 0.87076 \\
1.14 & 0.87286 \\
1.15 & 0.87493 \\
1.16 & 0.87698 \\
1.17 & 0.87900 \\
1.18 & 0.88100 \\
1.19 & 0.88298 \\
1.20 & 0.88493 \\
1.21 & 0.88686 \\
1.22 & 0.88877 \\
1.23 & 0.89065 \\
1.24 & 0.89251 \\
1.25 & 0.89435 \\
1.26 & 0.89617 \\
1.27 & 0.89796 \\
1.28 & 0.89973 \\
1.29 & 0.90147 \\
1.30 & 0.90320 \\
1.31 & 0.90490 \\
1.32 & 0.90658 \\
1.33 & 0.90824 \\
1.34 & 0.90988 \\
1.35 & 0.91149 \\
1.36 & 0.91309 \\
1.37 & 0.91466 \\
1.38 & 0.91621 \\
1.39 & 0.91774 \\
1.40 & 0.91924 \\
1.41 & 0.92073 \\
1.42 & 0.92220 \\
1.43 & 0.92364 \\
1.44 & 0.92507 \\
1.45 & 0.92647 \\
1.46 & 0.92785 \\
1.47 & 0.92922 \\
1.48 & 0.93056 \\
1.49 & 0.93189 \\
1.50 & 0.93319 \\
1.51 & 0.93448 \\
1.52 & 0.93574 \\
1.53 & 0.93699 \\
1.54 & 0.93822 \\
1.55 & 0.93943 \\
1.56 & 0.94062 \\
1.57 & 0.94179 \\
1.58 & 0.94295 \\
1.59 & 0.94408 \\
1.60 & 0.94520 \\
1.61 & 0.94630 \\
1.62 & 0.94738 \\
1.63 & 0.94845 \\
1.64 & 0.94950 \\
1.65 & 0.95053 \\
1.66 & 0.95154 \\
1.67 & 0.95254 \\
1.68 & 0.95352 \\
1.69 & 0.95449 \\
1.70 & 0.95543 \\
1.71 & 0.95637 \\
1.72 & 0.95728 \\
1.73 & 0.95818 \\
1.74 & 0.95907 \\
1.75 & 0.95994 \\
1.76 & 0.96080 \\
1.77 & 0.96164 \\
1.78 & 0.96246 \\
1.79 & 0.96327 \\
1.80 & 0.96407 \\
1.81 & 0.96485 \\
1.82 & 0.96562 \\
1.83 & 0.96638 \\
1.84 & 0.96712 \\
1.85 & 0.96784 \\
1.86 & 0.96856 \\
1.87 & 0.96926 \\
1.88 & 0.96995 \\
1.89 & 0.97062 \\
1.90 & 0.97128 \\
1.91 & 0.97193 \\
1.92 & 0.97257 \\
1.93 & 0.97320 \\
1.94 & 0.97381 \\
1.95 & 0.97441 \\
1.96 & 0.97500 \\
1.97 & 0.97558 \\
1.98 & 0.97615 \\
1.99 & 0.97670 \\
2.00 & 0.97725 \\
2.01 & 0.97778 \\
2.02 & 0.97831 \\
2.03 & 0.97882 \\
2.04 & 0.97932 \\
2.05 & 0.97982 \\
2.06 & 0.98030 \\
2.07 & 0.98077 \\
2.08 & 0.98124 \\
2.09 & 0.98169 \\
2.10 & 0.98214 \\
2.11 & 0.98257 \\
2.12 & 0.98300 \\
2.13 & 0.98341 \\
2.14 & 0.98382 \\
2.15 & 0.98422 \\
2.16 & 0.98461 \\
2.17 & 0.98500 \\
2.18 & 0.98537 \\
2.19 & 0.98574 \\
2.20 & 0.98610 \\
2.21 & 0.98645 \\
2.22 & 0.98679 \\
2.23 & 0.98713 \\
2.24 & 0.98745 \\
2.25 & 0.98778 \\
2.26 & 0.98809 \\
2.27 & 0.98840 \\
2.28 & 0.98870 \\
2.29 & 0.98899 \\
2.30 & 0.98928 \\
2.31 & 0.98956 \\
2.32 & 0.98983 \\
2.33 & 0.99010 \\
2.34 & 0.99036 \\
2.35 & 0.99061 \\
2.36 & 0.99086 \\
2.37 & 0.99111 \\
2.38 & 0.99134 \\
2.39 & 0.99158 \\
2.40 & 0.99180 \\
2.41 & 0.99202 \\
2.42 & 0.99224 \\
2.43 & 0.99245 \\
2.44 & 0.99266 \\
2.45 & 0.99286 \\
2.46 & 0.99305 \\
2.47 & 0.99324 \\
2.48 & 0.99343 \\
2.49 & 0.99361 \\
2.50 & 0.99379 \\
2.51 & 0.99396 \\
2.52 & 0.99413 \\
2.53 & 0.99430 \\
2.54 & 0.99446 \\
2.55 & 0.99461 \\
2.56 & 0.99477 \\
2.57 & 0.99492 \\
2.58 & 0.99506 \\
2.59 & 0.99520 \\
2.60 & 0.99534 \\
2.61 & 0.99547 \\
2.62 & 0.99560 \\
2.63 & 0.99573 \\
2.64 & 0.99585 \\
2.65 & 0.99598 \\
2.66 & 0.99609 \\
2.67 & 0.99621 \\
2.68 & 0.99632 \\
2.69 & 0.99643 \\
2.70 & 0.99653 \\
2.71 & 0.99664 \\
2.72 & 0.99674 \\
2.73 & 0.99683 \\
2.74 & 0.99693 \\
2.75 & 0.99702 \\
2.76 & 0.99711 \\
2.77 & 0.99720 \\
2.78 & 0.99728 \\
2.79 & 0.99736 \\
2.80 & 0.99744 \\
2.81 & 0.99752 \\
2.82 & 0.99760 \\
2.83 & 0.99767 \\
2.84 & 0.99774 \\
2.85 & 0.99781 \\
2.86 & 0.99788 \\
2.87 & 0.99795 \\
2.88 & 0.99801 \\
2.89 & 0.99807 \\
2.90 & 0.99813 \\
2.91 & 0.99819 \\
2.92 & 0.99825 \\
2.93 & 0.99831 \\
2.94 & 0.99836 \\
2.95 & 0.99841 \\
2.96 & 0.99846 \\
2.97 & 0.99851 \\
2.98 & 0.99856 \\
2.99 & 0.99861 \\
3.00 & 0.99865 \\
\end{longtable}



\end{document}
